% Kapitel: Das FourierGrating-Plugin (Beugungsgitter, RCWA / Fourier-Modal-Methode)
% Quellen: python/plugins/fouriergrating/ (fg_*.py, plugin.json, README.md), CHANGELOG.md,
% docs/adr/0010-multi-plugin-workflow.md
\chapter{Beugungsgitter mit FourierGrating}\label{kap:fouriergrating}
\index{FourierGrating}\index{Beugungsgitter}\index{RCWA}

Das Plugin \plugin{FourierGrating} berechnet die Beugung von Licht an periodischen Strukturen:
Liniengitter, Kreuzgitter, Schichtstapel mit Lamellen, Profilen oder beliebig geformten
Einschlüssen. Es verbindet \produkt{} mit dem Rechenkern \emph{FourierGrating}, einer
Implementierung der rigorosen Kopplungswellenanalyse (RCWA, auch Fourier-Modal-Methode genannt)
und der differentiellen Methode. Alles, was Sie für ein Gitter einstellen -- Struktur,
Materialien, Schichten, Beleuchtung --, steht im Panel \gui{Plugin}; die Ergebnisse erscheinen als
Tabellen, Balken, Kurven und Bilder in den Ergebnisfenstern der Anwendung.

Das Kapitel beschreibt das Plugin aus Anwendersicht: zuerst das Gittermodell mit seinen
Koordinaten, dann Materialien, Schichten und Beleuchtung, anschließend die Rechnungen (Einzelrechnung,
Scan, Aufwandsschätzung, Optimierung), das Nahfeld, Schichten aus CAD-Körpern und das Gitter als Block
im Node-Workflow. Den Abschluss bilden Arbeitsabläufe Schritt für Schritt, typische Fehler und die
Grenzen des Plugins. Die vollständige Liste aller Parameter mit Standardwerten und Wertebereichen
steht im Anhang, \autoref{kap:plugin-fouriergrating}.

\section{Überblick}\label{sec:fg-ueberblick}

\subsection{Was das Plugin kann}

\begin{itemize}
  \item \textbf{Einzelrechnung} (\gui{Solve}): die Beugungseffizienz jeder reflektierten und
    transmittierten Ordnung, die Summen $\Sigma R$ und $\Sigma T$, die Absorption und die
    Energiebilanz.
  \item \textbf{Scan} (\gui{Scan}): dieselben Größen über Listen oder Bereiche der Wellenlänge,
    der Einfallswinkel, der Polarisation und über einen beliebigen Zahlenwert der Konfiguration
    (Sweep), als Kurven.
  \item \textbf{Aufwandsschätzung} (\gui{Estimate}): Zahl der Rechnungen, Matrixgröße und grober
    Speicherbedarf, bevor Sie eine große Rechnung starten.
  \item \textbf{Optimierung} (\gui{Optimise}): Schichtdicken, Lamellenkanten, Perioden oder andere
    Zahlen der Konfiguration so wählen, dass eine gewichtete Summe von Effizienzen maximal oder
    minimal wird -- mit den exakten Gradienten von FourierGrating.
  \item \textbf{Littrow-Winkel} (\gui{Set Littrow angle}): den Einfallswinkel, unter dem eine
    reflektierte Ordnung in die Einfallsrichtung zurückläuft.
  \item \textbf{Nahfeld} (\gui{Near field}): das elektromagnetische Feld im und am Gitter als
    Bilder und Kurven.
  \item \textbf{Schichten aus CAD-Körpern} (\gui{Layer from CAD bodies}, \gui{Check CAD layers}):
    eine im CAD gezeichnete Elementarzelle als Gitterschicht.
  \item \textbf{Konfigurationsdateien} (\gui{Load configuration}, \gui{Save configuration}): das
    JSON-Format von FourierGrating lesen und schreiben, verlustfrei.
  \item \textbf{Blöcke für den Node-Workflow}: das Gitter \block{fg.grating} und den Schichtstapel
    \block{fg.stack}, die Strahlen und Felder im optischen Aufbau beugen bzw. filtern.
\end{itemize}

\subsection{Voraussetzungen}\label{sec:fg-voraussetzungen}
\index{fouriergrating (Python-Paket)}

Das Plugin wird mit der Anwendung installiert (Plugin-Kennung \code{fg}, Name
\plugin{FourierGrating}). Es braucht in der Python-Umgebung der Plugins die Pakete
\code{numpy}, \code{scipy} und \code{fouriergrating} in Version 0.18 oder neuer. Der Knopf
\gui{Create environment…} in \menu{Plugins,Settings…} installiert sie (siehe
\autoref{kap:installation} und \autoref{kap:plugins}). Fehlt das Paket \code{fouriergrating}, meldet
jede rechnende Aktion:

\begin{lstlisting}
The FourierGrating plugin needs the Python package fouriergrating in the plugin environment
(pip install fouriergrating, or the wheel of the FourierGrating repo).
\end{lstlisting}

Laden und Speichern von Konfigurationen, das Bearbeiten im Panel, \gui{Estimate} und
\gui{Set Littrow angle} funktionieren auch ohne das Paket. Zwei Funktionen haben weitere
Voraussetzungen:

\begin{itemize}
  \item \gui{Layer from CAD bodies} und \gui{Check CAD layers} brauchen ein \code{fouriergrating}
    mit CAD-Slicer (ein Build mit \code{FG\_ENABLE\_NURBFORGE}, siehe die Datei \datei{docs/cad.md}
    von FourierGrating). Ohne ihn wird die Schicht zwar angelegt, aber nicht geprüft.
  \item Die Workflow-Blöcke \block{fg.grating} und \block{fg.stack} stellt das Plugin nur bereit,
    wenn das Paket \code{rtt-fusion} mit Fernblöcken installiert ist (Modul
    \code{rtt\_fusion.graph.remote}).
\end{itemize}

\subsection{Wo Sie das Plugin finden}

\begin{itemize}
  \item Im Menü \menu{Plugins,FourierGrating} stehen alle Aktionen (alphabetisch nach ihrem
    internen Namen geordnet).
  \item In der Werkzeugspalte bildet das Plugin eine eigene Gruppe \gui{FourierGrating} mit einem
    Knopf je Aktion.
  \item Im Dock \gui{Plugin} wählen Sie oben \gui{FourierGrating}; darunter stehen die Parameter,
    unten die Knöpfe der Aktionen und \gui{Reset parameters} (setzt alle Parameter auf ihre
    Standardwerte zurück).
  \item Zwei Aktionen gelten einem einzelnen Block des Node-Workflows. Sie stehen nicht im Menü,
    sondern im Kontextmenü eines Gitter- oder Stapelblocks im Node-Editor
    (\autoref{sec:fg-block}).
\end{itemize}

Die Parameter werden mit dem Dokument gespeichert (pro Plugin). Ein geänderter Parameter markiert
das Dokument als geändert, ist aber kein Rückgängig-Schritt.

\subsection{Aufbau des Panels}\label{sec:fg-panel}

Das Panel ist in Gruppen gegliedert. Ganz oben stehen ohne eigenen Rahmen die beiden Listen
\gui{Materials} und \gui{Layers (bottom to top)}, darunter die Gruppenkästen in dieser
Reihenfolge:

\begin{referenz}{>{\raggedright\arraybackslash}p{3.2cm}>{\raggedright\arraybackslash}X}
\toprule
Gruppe & Inhalt \\
\midrule
\endhead
(ohne Gruppe) & \gui{Materials} (Liste der Materialien), \gui{Layers (bottom to top)} (Schichtstapel) \\
\gui{File} & \gui{Configuration} (Datei zum Laden), \gui{Save as} (Datei zum Speichern) \\
\gui{Structure} & \gui{Dimension}, \gui{Solver}, \gui{Orders Nx}, \gui{Orders Nz}, \gui{Period x},
  \gui{Period z}, \gui{Superstrate}, \gui{Substrate}, \gui{Solve: also write CSV} \\
\gui{Incidence} & \gui{$\lambda$ / \si{\micro\metre}}, \gui{$\varphi$ / deg}, \gui{$\psi$ / deg},
  \gui{$\delta$ / deg}, \gui{Littrow order} \\
\gui{Scan} & \gui{x axis}, \gui{Orders from efficiency}, \gui{At most curves}, \gui{Also write CSV} \\
\gui{Near field} & \gui{Show}, \gui{y from}, \gui{y to}, \gui{Horizontal cut at y},
  \gui{Vertical cut at z (2D)}, \gui{Points x}, \gui{Points y}, \gui{Points z (2D)}, \gui{Periods},
  \gui{Also write .npz} \\
\gui{Optimisation} & \gui{Goal}, \gui{Objective L = $\Sigma$ weight $\cdot$ efficiency},
  \gui{Parameters}, \gui{Iterations}, \gui{Apply the optimum to the panel} \\
\gui{CAD layer} & \gui{Bodies}, \gui{Bodies drawn in}, \gui{STEP file}, \gui{Around the bodies},
  \gui{RCWA slices per band}, \gui{Preview at height}, \gui{Preview SVG} \\
\gui{Advanced} & \gui{Further settings (JSON)}: alle übrigen Schlüssel der Konfiguration \\
\bottomrule
\end{referenz}

Fahren Sie mit der Maus über ein Feld, zeigt ein Tooltip seine Kurzbeschreibung. Ungültige
Eingaben (etwa fehlerhaftes JSON oder eine falsche Tabellenzelle) werden rot markiert und sperren
die Aktionen, bis sie behoben sind.

\subsection{Wo die Ergebnisse erscheinen}

Jede rechnende Aktion legt im Dock \gui{Results} einen neuen Ergebnissatz an (etwa
\enquote{FourierGrating $\lambda$ = 0.6328 \si{\micro\metre}}). Ergebnissätze lassen sich
ein- und ausblenden, vergleichen und löschen; sie werden mit dem Dokument gespeichert. Die
einzelnen Ergebnisse landen je nach Art in

\begin{itemize}
  \item dem Fenster \gui{Table} (Ordnungstabelle, Parameter der Optimierung, Bänder und Befunde der
    CAD-Prüfung),
  \item dem Fenster \gui{Plot} (Balken der Effizienzen, Kurven des Scans und der Optimierung,
    horizontaler Nahfeldschnitt eines 1D-Gitters),
  \item dem Fenster \gui{Image} (Nahfeldschnitte als Farbbilder),
  \item dem Dock \gui{Plugin log} (Summen, Energiebilanz, Warnungen, Schritte der Optimierung) und
  \item der Statusleiste (eine Zusammenfassung nach dem Ende der Aktion).
\end{itemize}

Ein Fehler, den Sie selbst beheben können -- eine fehlende Datei, ein unbekanntes Material --,
erscheint als Meldungsfenster mit dem Text des Plugins. Alle anderen Fehler stehen mit
Einzelheiten im \gui{Plugin log}. Eine laufende Aktion brechen Sie mit dem Knopf \gui{Cancel} in der
Statusleiste ab.

\begin{hinweis}
Eine einzelne Rechnung von FourierGrating, die bereits läuft, lässt sich nicht mittendrin
unterbrechen. \gui{Optimise} prüft den Abbruch zwischen zwei Auswertungen; bei \gui{Solve},
\gui{Scan} und \gui{Near field} wirkt \gui{Cancel} erst, wenn die Rechnung zurückkehrt -- oder die
Anwendung beendet den Plugin-Prozess nach einigen Sekunden. Schätzen Sie große Rechnungen deshalb
vorher mit \gui{Estimate} ab (\autoref{sec:fg-estimate}).
\end{hinweis}

% ---------------------------------------------------------------------------------------------
\section{Das Gittermodell}\label{sec:fg-modell}
\index{Gittermodell}

\subsection{Koordinaten und Schichtung}\label{sec:fg-koordinaten}

FourierGrating beschreibt ein Gitter als Stapel von Schichten zwischen zwei Halbräumen:

\begin{itemize}
  \item Die \textbf{y-Achse} ist die Stapelrichtung. Unten liegt das \textbf{Substrat}
    ($y < 0$), oben das \textbf{Superstrat}. Die Schichten liegen dazwischen, die unterste beginnt
    bei $y = 0$.
  \item Das Licht fällt \textbf{von oben}, aus dem Superstrat, ein. Reflektierte Ordnungen laufen
    zurück ins Superstrat (Seite \code{R}), transmittierte ins Substrat (Seite \code{T}).
  \item Die Struktur ist periodisch in \textbf{x} mit der Periode $\Lambda_x$ (\gui{Period x}),
    bei einem Kreuzgitter (2D) zusätzlich in \textbf{z} mit $\Lambda_z$ (\gui{Period z}).
  \item Alle Längen sind in Mikrometern (\si{\micro\metre}), alle Winkel in Grad.
\end{itemize}

Diese Konvention gilt im ganzen Plugin: im Panel, in den Konfigurationsdateien, beim Nahfeld
(\gui{y from}/\gui{y to} zählen ab der Unterkante des Stapels) und bei den CAD-Schichten (dort
wird die CAD-Achse $+Z$ zur Stapelrichtung $y$, siehe \autoref{sec:fg-cad}).

\subsection{Dimension: Liniengitter und Kreuzgitter}

\gui{Dimension} wählt zwischen
\begin{description}
  \item[\code{1D}] einem Liniengitter: periodisch nur in x, unendlich ausgedehnt und konstant in z.
    Die Ordnungen tragen einen Index $n$.
  \item[\code{2D}] einem Kreuzgitter: periodisch in x und z. Die Ordnungen tragen zwei Indizes
    $(n, m)$; \gui{Period z} muss dann größer als null sein.
\end{description}

Lamellen- und isotrope Profilschichten eines Liniengitters sind auch in einem 2D-Stapel erlaubt;
sie sind dann in z konstant.

\subsection{Fourier-Ordnungen und Matrixgröße}\label{sec:fg-ordnungen}
\index{Fourier-Ordnungen}\index{Abbruch (Truncation)}

Die Methode entwickelt Feld und Material in Fourier-Reihen. \gui{Orders Nx} und \gui{Orders Nz}
legen fest, wie viele Glieder mitgenommen werden: von $-N_x$ bis $+N_x$ in x und von $-N_z$ bis
$+N_z$ in z. Die Zahl der Fourier-Glieder ist
\[
  n_f = 2N_x + 1 \quad\text{(1D)}, \qquad n_f = (2N_x + 1)(2N_z + 1) \quad\text{(2D)},
\]
und die Matrix des Eigenwertproblems hat die Größe $n = 4 n_f$. Rechenzeit und Speicher wachsen
mit $n^3$ bzw. $n^2$ -- ein 2D-Gitter mit $N_x = N_z = 10$ hat $n = 1764$, ein 1D-Gitter mit
$N_x = 10$ nur $n = 84$.

\begin{tipp}
Prüfen Sie die Konvergenz: rechnen Sie mit $N_x$, dann mit einem deutlich größeren Wert (etwa dem
Doppelten). Ändern sich die Effizienzen, die Sie interessieren, nicht mehr in den benötigten
Stellen, reicht der kleinere Wert. Metallische Lamellen in TM-Polarisation konvergieren langsam und
brauchen viele Ordnungen.
\end{tipp}

Die Standardwerte sind $N_x = 10$ und $N_z = 0$, der erlaubte Bereich 0 bis 200.

\subsection{Rechenverfahren}\label{sec:fg-solver}
\index{differentielle Methode}

\gui{Solver} wählt das Verfahren:
\begin{description}
  \item[\code{rcwa}] die RCWA (Fourier-Modal-Methode): jede Schicht wird in y in Scheiben
    konstanter Struktur zerlegt, in jeder Scheibe wird ein Eigenwertproblem gelöst. Für
    Lamellen- und homogene Schichten ist das exakt; Profile werden durch Scheiben angenähert.
    Die Scheibenzahl regelt der Schlüssel \code{rcwa} unter \gui{Advanced}
    (\code{\{\textquotedbl{}rcwa\textquotedbl{}: \{\textquotedbl{}slices\textquotedbl{}: 60\}\}}) bzw. \code{slices} der einzelnen Profilschicht.
  \item[\code{rk}] die differentielle Methode: das Feld wird mit einem impliziten
    Runge-Kutta-Verfahren durch die Schicht integriert. Sie eignet sich für stetig veränderliche
    Profile. Einstellungen (Tableau, Toleranz) stehen ebenfalls unter \gui{Advanced}
    (Schlüssel \code{rk}).
\end{description}

\subsection{Superstrat und Substrat}

\gui{Superstrate} und \gui{Substrate} nennen die Materialien der beiden Halbräume -- über ihren
Namen in der Materialliste (Standard: \code{air} und \code{glass}). Ein Name, der in der Liste
fehlt, führt beim Rechnen zu einer Meldung von FourierGrating.

\begin{warnung}
Das Superstrat muss verlustfrei sein (reelle Brechzahl), denn aus ihm fällt das Licht ein.
FourierGrating prüft das beim Laden der Konfiguration.
\end{warnung}

\subsection{Gittergleichung und Ordnungen}\label{sec:fg-gittergleichung}
\index{Gittergleichung}

Für eine ebene Welle aus dem Superstrat mit der Brechzahl $\nu$ unter dem Polarwinkel $\varphi$
(gemessen gegen die y-Achse, also gegen die Flächennormale) und dem Azimut $\psi$ (gemessen gegen
x) lautet die x-Komponente des normierten Wellenvektors der Ordnung $n$
\[
  \alpha_n = \nu \sin\varphi \cos\psi + n\,\frac{\lambda}{\Lambda_x},
\]
in 2D zusätzlich $\gamma_m = \nu \sin\varphi \sin\psi + m\,\lambda/\Lambda_z$. Eine Ordnung
breitet sich in einem Medium mit der Brechzahl $n_\text{M}$ nur aus, wenn
$\alpha_n^2 + \gamma_m^2 < n_\text{M}^2$; sonst ist sie evaneszent und trägt keine Leistung fort.
Deshalb hängt die Zahl der Zeilen in der Ordnungstabelle von $\lambda$, $\Lambda$ und den Winkeln
ab, und im Scan kann eine Ordnung an manchen Punkten fehlen.

% ---------------------------------------------------------------------------------------------
\section{Materialien}\label{sec:fg-materialien}
\index{Material (FourierGrating)}\index{Brechzahl}

Die Liste \gui{Materials} enthält alle Materialien, auf die sich Superstrat, Substrat und Schichten
beziehen. Standardmäßig stehen dort \code{air} ($n = 1$) und \code{glass} ($n = 1{,}5$).

\subsection{Materialien bearbeiten}

Die Liste hat die Knöpfe \gui{Add} (neuer Eintrag), \gui{Edit…} (öffnet den Eintrag in einem
Dialog \gui{Edit entry}), \gui{Remove}, $\uparrow$ und $\downarrow$ (Reihenfolge ändern). Ein
Eintrag hat diese Felder:

\begin{description}
  \item[\gui{Name}] der Name, unter dem Schichten, Superstrat und Substrat das Material nennen.
    Groß- und Kleinschreibung zählt.
  \item[\gui{Given as}] \code{n} -- die komplexe Brechzahl $n + \mathrm{i}k$ aus den Feldern
    \gui{n} und \gui{k}; \code{json} -- jede andere Form, die FourierGrating kennt, als JSON im
    Feld \gui{Material as JSON (json)}.
  \item[\gui{n}, \gui{k}] Real- und Imaginärteil der Brechzahl. $k > 0$ bedeutet Absorption.
    Ist $k = 0$, wird das Material als reelle Zahl gespeichert, sonst als Paar \code{[n, k]}.
\end{description}

\subsection{Weitere Materialformen (JSON)}

Mit \gui{Given as} = \code{json} geben Sie das Material so ein, wie es in der Konfigurationsdatei
unter seinem Namen stünde. FourierGrating kennt unter anderem (Format nach seiner Dokumentation
\datei{docs/config\_format.md}):

\begin{itemize}
  \item einen vollständigen Permittivitätstensor (\code{eps\_tensor}) für anisotrope Materialien,
  \item Hauptachsenwerte mit Euler-Winkeln (\code{eps\_principal}, \code{euler\_deg}),
  \item dispersive Modelle, die an jeder Wellenlänge eines Scans ausgewertet werden: Sellmeier,
    Cauchy, Drude-Lorentz und Tabellen.
\end{itemize}

Ein Beispiel: Quarzglas nach Sellmeier.

\begin{lstlisting}
{ "n": { "sellmeier": { "A": 1.0, "B": [0.6961663, 0.4079426, 0.8974794],
                        "C_um2": [0.00467914826, 0.0135120631, 97.9340025] } } }
\end{lstlisting}

Materialien, die Sie als JSON eingeben, reicht das Plugin unverändert an FourierGrating weiter; beim
Laden einer Datei landen alle Materialien, die keine einfache Brechzahl sind, automatisch in dieser
Form im Panel.

\begin{hinweis}
\gui{Set Littrow angle} braucht die Brechzahl des Superstrats als Zahl. Ist das Superstrat als JSON
gegeben, meldet die Aktion \enquote{the superstrate … needs a refractive index n (not JSON)}.
\end{hinweis}

% ---------------------------------------------------------------------------------------------
\section{Der Schichtstapel}\label{sec:fg-schichten}
\index{Schichtstapel}\index{Lamellengitter}

Die Liste \gui{Layers (bottom to top)} beschreibt die Schichten zwischen Substrat und Superstrat --
\textbf{von unten nach oben}: der erste Eintrag liegt direkt auf dem Substrat, der letzte grenzt an
das Superstrat. Das ist auch die Zählung der JSON-Zeiger: \code{/layers/0} ist die unterste
Schicht. Mit $\uparrow$ und $\downarrow$ verschieben Sie eine Schicht im Stapel. Ein leerer Stapel
ist erlaubt: dann rechnet FourierGrating die ebene Grenzfläche zwischen Superstrat und Substrat.

Jeder Eintrag hat einen \gui{Type}; die übrigen Felder gelten je nach Typ.

\subsection{Homogene Schicht}

\gui{Type} = \code{homogeneous}: eine ebene Schicht der Dicke \gui{Thickness} (in
\si{\micro\metre}) aus dem Material \gui{Material (homogeneous)}. Damit bauen Sie Beschichtungen,
Zwischenschichten und Deckschichten.

\subsection{Lamellenschicht}\label{sec:fg-lamellen}

\gui{Type} = \code{lamellar1d}: eine Schicht der Dicke \gui{Thickness}, die in x aus Lamellen
verschiedener Materialien besteht -- das klassische binäre Gitter. Die Lamellen stehen in der
Tabelle \gui{Lamellae (lamellar1d)} mit den Spalten \gui{from x} (Beginn der Lamelle in
\si{\micro\metre}) und \gui{material}. Jede Lamelle reicht von ihrem \gui{from x} bis zum
\gui{from x} der nächsten Zeile, die letzte bis \gui{End (lamellar1d)}. Zeilen fügen Sie mit
\gui{Add row} an und löschen sie mit \gui{Remove row}.

Ein Gitter mit Periode \qty{1}{\micro\metre}, dessen erste Hälfte Glas und zweite Hälfte Luft ist:

\begin{center}
\begin{tabular}{rl}
\toprule
\gui{from x} & \gui{material} \\
\midrule
0.0 & glass \\
0.5 & air \\
\bottomrule
\end{tabular}
\qquad \gui{End (lamellar1d)} = 1.0
\end{center}

In der Konfigurationsdatei wird daraus \code{\textquotedbl{}breakpoints\_um\textquotedbl{}: [0.0, 0.5, 1.0]} und
\code{\textquotedbl{}materials\textquotedbl{}: [\textquotedbl{}glass\textquotedbl{}, \textquotedbl{}air\textquotedbl{}]}.

\begin{warnung}
Die Kanten müssen bei 0 beginnen und genau bei der Periode $\Lambda_x$ enden: die erste Zeile hat
\gui{from x} = 0, \gui{End} ist gleich \gui{Period x}. FourierGrating prüft das
\textbf{nicht}. Ändern Sie die Periode, ohne \gui{End} anzupassen, ist das Ergebnis ohne Meldung
falsch.
\end{warnung}

\subsection{Profile und weitere Schichttypen (JSON)}\label{sec:fg-profile}
\index{Profilgitter}

\gui{Type} = \code{json}: jede andere Schicht von FourierGrating, eingegeben im Feld
\gui{Layer as JSON (type json)} genau so, wie sie in der Konfigurationsdatei unter
\code{layers} stünde. Laut der Formatbeschreibung von FourierGrating gibt es unter anderem:

\begin{description}
  \item[\code{profile1d}] ein stetiges Oberflächenprofil eines Liniengitters zwischen zwei
    Materialien (\code{below}, \code{above}) mit der Tiefe \code{depth\_um}. Formen:
    \code{sinus}, \code{triangle} (Dreieck mit der Spitze bei \code{apex\_x\_um}) und
    \code{points} (Stützstellen \code{x\_um}, \code{y\_um}, stückweise linear).
  \item[\code{profile2d}] ein Profil eines Kreuzgitters (\code{egg\_crate}, \code{sinus\_x}).
  \item[\code{polygon2d}] eine Schicht eines Kreuzgitters mit polygonalen oder rechteckigen
    Einschlüssen in einem Hintergrundmaterial.
  \item[\code{contour2d}] Einschlüsse mit gekrümmten Rändern: Kreise, Ellipsen, Polygone,
    NURBS-Konturen, auch mit Löchern.
  \item[\code{cad}] Körper aus einer STEP-Datei, die FourierGrating in Bänder zerlegt -- diese
    Schicht legt das Plugin für Sie an (\autoref{sec:fg-cad}).
\end{description}

Ein sinusförmiges Reliefgitter in Aluminium unter Luft, \qty{0.2}{\micro\metre} tief:

\begin{lstlisting}
{ "type": "profile1d", "shape": "sinus", "depth_um": 0.2,
  "below": "al", "above": "air", "slices": 0, "factorization": "auto" }
\end{lstlisting}

Das Plugin prüft den Inhalt eines JSON-Eintrags nicht selbst; Fehler meldet FourierGrating beim
Rechnen (\enquote{FourierGrating: …}). Ungültiges JSON markiert das Feld rot.

\begin{tipp}
Wenn Sie eine Konfigurationsdatei mit solchen Schichten laden, zeigt das Panel sie als Einträge vom
Typ \code{json}; nur homogene und Lamellenschichten in genau der einfachen Form erscheinen als
Felder. Beim Speichern wird jede Schicht wieder so geschrieben, wie sie war.
\end{tipp}

% ---------------------------------------------------------------------------------------------
\section{Beleuchtung und Polarisation}\label{sec:fg-beleuchtung}
\index{Polarisation}\index{Einfallswinkel}

Die Gruppe \gui{Incidence} beschreibt die einfallende ebene Welle:

\begin{description}
  \item[\gui{$\lambda$ / \si{\micro\metre}}] die Vakuumwellenlänge in Mikrometern
    (Standard 0.6328, die Linie des HeNe-Lasers).
  \item[\gui{$\varphi$ / deg}] der Polarwinkel des Einfalls gegen die Flächennormale (die
    y-Achse). $\varphi = 0$ ist senkrechter Einfall.
  \item[\gui{$\psi$ / deg}] der Azimut der Einfallsebene gegen die x-Achse. $\psi = 0$: die
    Einfallsebene steht senkrecht auf den Gitterlinien eines Liniengitters (klassische
    Anordnung); $\psi \neq 0$: konischer Einfall.
  \item[\gui{$\delta$ / deg}] der Polarisationswinkel: \code{0} ist s-Polarisation (TE),
    \code{90} p-Polarisation (TM).
\end{description}

\subsection{Zahl, Liste oder Bereich}\label{sec:fg-scan-syntax}

Jedes der vier Felder nimmt eine von drei Formen an:

\begin{center}
\begin{tabularx}{\linewidth}{ll>{\raggedright\arraybackslash}X}
\toprule
Form & Eingabe & Bedeutung \\
\midrule
Zahl & \code{0.6328} & ein Wert \\
Liste & \code{0.5, 0.6, 0.7} & die genannten Werte \\
Bereich & \code{0.5:0.7:21} & von 0,5 bis 0,7 in 21 gleichen Schritten (beide Enden eingeschlossen) \\
\bottomrule
\end{tabularx}
\end{center}

Schreiben Sie Dezimalzahlen in diesen Feldern immer mit Punkt. Eine Liste trennen Sie mit
Kommas, einen Bereich mit Doppelpunkten in der Reihenfolge \code{von:bis:Schritte}.

Mehrere Listen und Bereiche ergeben einen Scan über das \textbf{kartesische Produkt}: drei
Wellenlängen und zwei Polarisationen sind sechs Rechenpunkte. \gui{Scan} rechnet alle Punkte;
\gui{Solve}, \gui{Near field} und \gui{Optimise} rechnen nur den \textbf{ersten} Punkt (die erste
Wellenlänge, den ersten Winkel, die erste Polarisation).

\begin{hinweis}
Ein leeres Feld lässt den Schlüssel in der Konfiguration weg. FourierGrating setzt dann seinen
eigenen Standard ein. Die Wellenlänge sollten Sie immer angeben.
\end{hinweis}

\section{Littrow-Anordnung}\label{sec:fg-littrow}
\index{Littrow-Anordnung}

In der Littrow-Anordnung läuft eine reflektierte Beugungsordnung genau in die Einfallsrichtung
zurück -- typisch für Gitter in Laserresonatoren und Spektrometern. Für die Ordnung $n$ gilt dann
$\alpha_n = -\alpha_0$, also
\[
  \sin\varphi = -\frac{n\,\lambda}{2\,\Lambda_x\,\nu\,\cos\psi}.
\]
Dabei ist $\nu$ die Brechzahl des Superstrats, $\lambda$ die erste Wellenlänge und $\psi$ der erste
Azimut des Panels.

\begin{schritte}
  \item Tragen Sie in der Gruppe \gui{Incidence} unter \gui{Littrow order} die gewünschte Ordnung
    ein (Standard $-1$, erlaubt $-50$ bis $50$).
  \item Wählen Sie \menu{Plugins,FourierGrating,Set Littrow angle} oder den gleichnamigen Knopf.
  \item Das Plugin schreibt den Winkel in das Feld \gui{$\varphi$ / deg} (mit zehn gültigen
    Stellen) und meldet ihn in der Statusleiste, etwa \enquote{$\varphi$ = 18.4454\si{\degree} (Littrow, order -1)}.
\end{schritte}

Beispiel: $\lambda = \qty{0.6328}{\micro\metre}$, $\Lambda_x = \qty{1}{\micro\metre}$, Luft
($\nu = 1$), $\psi = 0$, Ordnung $-1$: $\sin\varphi = 0{,}3164$, also
$\varphi \approx \ang{18.45}$.

Die Aktion meldet einen Fehler, wenn \gui{Period x} nicht größer als null ist, wenn das Superstrat
nicht als Brechzahl gegeben ist, oder wenn die Ordnung keinen Littrow-Winkel hat (Ordnung 0 oder
$|\sin\varphi| > 1$, die Meldung nennt den Wert).

\begin{tipp}
\gui{Set Littrow angle} überschreibt das Feld \gui{$\varphi$ / deg} mit einer einzelnen Zahl; ein
vorher eingetragener Bereich ist danach weg.
\end{tipp}

% ---------------------------------------------------------------------------------------------
\section{Rechnen}\label{sec:fg-rechnen}

\subsection{Einzelrechnung: Solve}\label{sec:fg-solve}
\index{Beugungseffizienz}

\gui{Solve} rechnet die Konfiguration des Panels für den ersten Punkt der Beleuchtung und zeigt die
Effizienzen aller Ordnungen, die sich im Superstrat oder Substrat ausbreiten.

\begin{schritte}
  \item Stellen Sie Struktur, Materialien, Schichten und Beleuchtung ein.
  \item Wählen Sie \menu{Plugins,FourierGrating,Solve} oder den Knopf \gui{Solve} im Panel.
  \item Die Statusleiste zeigt \enquote{solving …}, nach dem Ende eine Zusammenfassung der Form (die Zahlen hier nur als Beispiel)
    \enquote{FourierGrating: $\Sigma$R 0.0400, $\Sigma$T 0.9600, A 0.0000 (8 orders)}.
\end{schritte}

Das Ergebnis besteht aus

\begin{description}
  \item[Tabelle \enquote{Orders}] im Fenster \gui{Table}, eine Zeile je Ordnung mit den Spalten
    \code{side} (\code{R} reflektiert, \code{T} transmittiert), \code{n}, \code{m} (in 1D immer 0),
    \code{efficiency}, \code{s part} und \code{p part}. Die Effizienz ist der Anteil der einfallenden
    Leistung, den die Ordnung fortträgt; \code{s part} und \code{p part} teilen ihn auf die beiden
    Polarisationen der gebeugten Welle auf (in isotropen Medien).
  \item[Balken \enquote{Efficiency per order}] im Fenster \gui{Plot}: je Ordnung ein Balken,
    beschriftet \code{R 0}, \code{R -1}, \code{T 1} usw. (in 2D \code{R 1,0}).
  \item[Protokollzeile] im \gui{Plugin log}: $\Sigma R$, $\Sigma T$, die Absorption $A$ und,
    wenn FourierGrating sie liefert, die Energiebilanz; davor jede Warnung von FourierGrating
    (\enquote{warning: …}).
\end{description}

Der Ergebnissatz heißt nach der Wellenlänge, etwa \enquote{FourierGrating $\lambda$ = 0.6328
\si{\micro\metre}}; bei einer Liste steht die ganze Liste im Namen, gerechnet wird trotzdem nur
der erste Punkt.

\paragraph{Energiebilanz.} Für ein verlustfreies Gitter (alle $k = 0$) muss
$\Sigma R + \Sigma T = 1$ sein, $A = 0$. Weicht die Summe merklich ab, ist die Rechnung nicht
konvergiert oder falsch eingestellt (zu wenige Ordnungen, zu wenige Scheiben, falsche
Lamellenkanten). Mit absorbierenden Materialien ist $A = 1 - \Sigma R - \Sigma T > 0$.

\paragraph{Als CSV schreiben.} Steht in \gui{Solve: also write CSV} (Gruppe \gui{Structure}) eine
Datei, schreibt \gui{Solve} die Ordnungstabelle zusätzlich dorthin -- mit den Spalten
\code{side, n, m, efficiency, eff\_s, eff\_p}. Das Protokoll meldet \enquote{\dots{} orders written to
\dots}. Lassen Sie das Feld leer, wenn Sie keine Datei wollen.

\subsection{Parameterscan: Scan}\label{sec:fg-scan}
\index{Scan (FourierGrating)}\index{Sweep}

\gui{Scan} rechnet \emph{alle} Punkte der Listen und Bereiche in der Gruppe \gui{Incidence} und
eines Sweeps -- parallel, in FourierGrating -- und zeigt die Effizienzen als Kurven.

\paragraph{Was gescannt wird.}
\begin{itemize}
  \item Die vier Felder der Beleuchtung: jede Liste oder jeder Bereich (\autoref{sec:fg-scan-syntax}).
    Mehrere variierende Größen ergeben das kartesische Produkt.
  \item Ein \textbf{Sweep} über eine beliebige Zahl der Konfiguration, etwa eine Schichtdicke.
    Ihn tragen Sie unter \gui{Advanced} in \gui{Further settings (JSON)} ein:
\begin{lstlisting}
{ "sweep": { "parameter": "/layers/0/thickness_um",
             "values": { "from": 0.1, "to": 0.3, "steps": 11 } } }
\end{lstlisting}
    \code{parameter} ist ein JSON-Zeiger in die Konfiguration (\autoref{sec:fg-zeiger}),
    \code{values} eine Zahl, eine Liste oder ein Bereich. Die Sweep-Werte laufen nacheinander, die
    Punkte jedes Werts parallel.
\end{itemize}

Gibt es nur einen einzigen Punkt, meldet \gui{Scan}: \enquote{the scan has a single point – give a
list or a range for $\lambda$, an angle or a sweep}.

\paragraph{Die Kurven.} Im Fenster \gui{Plot} erscheinen
\begin{itemize}
  \item $\Sigma R$ und $\Sigma T$ über der x-Achse,
  \item jede Ordnung, deren Effizienz irgendwo den Wert \gui{Orders from efficiency} erreicht
    (Standard 0,01), als eigene Kurve, benannt \code{R 0}, \code{T -1} (in 2D \code{R 1,0}).
\end{itemize}

Die x-Achse wählt \gui{x axis}: \code{auto} nimmt die Größe mit den meisten Werten (bei Gleichstand
zuerst $\lambda$); \code{lambda\_um}, \code{phi\_deg}, \code{psi\_deg}, \code{delta\_deg} oder
\code{sweep\_value} erzwingen eine bestimmte. Variiert die gewählte Größe nicht, meldet das Plugin,
welche Größen variieren. Jede \emph{weitere} variierende Größe erzeugt eigene Kurven; ihr Wert
steht im Kurvennamen, etwa \code{R 0 ($\delta$ = 90)} oder \code{$\Sigma$T ($\varphi$ = 10, sweep =
0.2)} -- beim Sweep steht dort der Name des Sweep-Parameters.

Eine Ordnung, die an manchen Punkten evaneszent ist (\autoref{sec:fg-gittergleichung}), hat dort
keine Punkte; ihre Kurve hat eine Lücke oder beginnt später.

\paragraph{Zu viele Kurven.} Mehr als \gui{At most curves} (Standard 24) Kurven zeigt das Plugin
nicht; es nimmt die ersten und schreibt eine Warnung ins Protokoll (\enquote{raise 'Orders from
efficiency' or 'At most curves'}). Erhöhen Sie \gui{Orders from efficiency}, um schwache Ordnungen
auszublenden, oder \gui{At most curves}.

\paragraph{Als CSV schreiben.} \gui{Also write CSV} in der Gruppe \gui{Scan} schreibt jede Ordnung an
jedem Punkt als Zeile -- im Format der Datei \datei{efficiencies.csv} von FourierGrating, komplexe
Spalten (Amplituden) aufgeteilt in \code{\_re} und \code{\_im}. Die Statusleiste nennt die Zahl der
Zeilen.

Die Zusammenfassung in der Statusleiste lautet etwa \enquote{21 points, 6 curves over $\lambda$ /
\si{\micro\metre}}.

\begin{tipp}
Für einen Vergleich von s- und p-Polarisation tragen Sie bei \gui{$\delta$ / deg} die Liste
\code{0, 90} ein und scannen über $\lambda$: Sie erhalten jede Kurve zweimal, mit
\code{($\delta$ = 0)} und \code{($\delta$ = 90)} im Namen.
\end{tipp}

\subsection{Aufwand abschätzen: Estimate}\label{sec:fg-estimate}
\index{Speicherbedarf}

\gui{Estimate} rechnet nichts, sondern schätzt aus der Konfiguration, wie groß eine Rechnung wird.
Die Statusleiste zeigt etwa

\begin{center}
\code{21 solve(s) $\cdot$ matrix size n = 84 $\cdot$ memory roughly 1.3 MiB (rule of thumb)}
\end{center}

\begin{itemize}
  \item \textbf{Rechnungen}: Zahl der Scanpunkte (Produkt der Längen der vier Beleuchtungsfelder) mal
    Zahl der Sweep-Werte.
  \item \textbf{Matrixgröße} $n = 4 n_f$ (\autoref{sec:fg-ordnungen}).
  \item \textbf{Speicher} nach einer Faustregel aus GUI-FourierGrating: gleichzeitig gerechnete
    Punkte $\times\, 12\, n^2 \times 16$ Byte. Gleichzeitig laufen höchstens so viele Punkte, wie
    der Rechner Kerne hat.
\end{itemize}

Übersteigt die Schätzung 70\,\% des Arbeitsspeichers, steht im Protokoll eine Warnung: \enquote{the
estimated memory is \dots{} of the RAM – fewer orders (Nx, Nz) or fewer parallel scan points
(threads.scan under Advanced)}. \gui{Scan} und \gui{Optimise} prüfen das vor dem Start ebenfalls und
warnen im Protokoll; sie rechnen trotzdem.

\begin{center}
\begin{tabular}{lrrr}
\toprule
Beispiel & $n_f$ & $n$ & Speicher je Punkt \\
\midrule
1D, $N_x = 10$ & 21 & 84 & \qty{1.3}{MiB} \\
1D, $N_x = 20$ & 41 & 164 & \qty{4.9}{MiB} \\
2D, $N_x = N_z = 10$ & 441 & 1764 & \qty{570}{MiB} \\
\bottomrule
\end{tabular}
\end{center}

Ein 2D-Scan mit $N_x = N_z = 10$ auf acht Kernen braucht nach dieser Regel also rund
\qty{4.5}{GiB}.

\begin{hinweis}
Die Schätzung ist eine Warnung, keine Vorhersage. Mit dem Runge-Kutta-Verfahren, vielen Scheiben
oder der GPU kann der wirkliche Bedarf stark abweichen. Die Zahl der parallelen Scanpunkte begrenzen
Sie unter \gui{Advanced} mit \code{\{\textquotedbl{}threads\textquotedbl{}: \{\textquotedbl{}scan\textquotedbl{}: 2\}\}}.
\end{hinweis}

\subsection{Optimieren: Optimise}\label{sec:fg-optimieren}
\index{Optimierung (Gitter)}\index{L-BFGS-B}

\gui{Optimise} sucht Werte für ausgewählte Zahlen der Konfiguration -- Schichtdicken,
Lamellenkanten, Perioden --, die eine Zielfunktion maximieren oder minimieren. Das Verfahren ist
L-BFGS-B (aus \code{scipy}) mit den exakten Gradienten, die FourierGrating über die adjungierte
Methode liefert: ein Optimierungsschritt kostet etwa 1,3 bis 2,5 Rechnungen, unabhängig von der Zahl
der Parameter. Gerechnet wird am ersten Punkt der Beleuchtung.

\paragraph{Die Zielfunktion.} Die Tabelle \gui{Objective L = $\Sigma$ weight $\cdot$ efficiency}
legt fest, welche Effizienzen zählen:
\[
  L = \sum_\text{Zeilen} w \cdot \eta(\text{side}, n, m).
\]
Jede Zeile hat die Spalten \gui{side} (\code{R} oder \code{T}), \gui{n}, \gui{m} (in 1D 0) und
\gui{weight} ($w$, Standard 1). Zeilen mit derselben Ordnung werden zusammengezählt. \gui{Goal} wählt
\code{maximise} (Standard) oder \code{minimise}.

Beispiele:
\begin{itemize}
  \item Maximale Effizienz der $-1$. reflektierten Ordnung: eine Zeile \code{R, -1, 0, 1}.
  \item Möglichst wenig Rückreflex: eine Zeile \code{R, 0, 0, 1}, \gui{Goal} = \code{minimise}.
  \item Eine Ordnung stärken und eine andere schwächen: \code{R, -1, 0, 1} und \code{R, 0, 0, -1},
    \gui{Goal} = \code{maximise}.
\end{itemize}

\paragraph{Die Parameter.}\label{sec:fg-zeiger}
\index{JSON-Zeiger}
Die Tabelle \gui{Parameters} nennt die Zahlen, die verändert werden dürfen, je mit Unter- und
Obergrenze. Die Spalte \gui{JSON pointer} adressiert die Zahl in der Konfiguration (Schreibweise nach
RFC 6901: mit \code{/} beginnend, Listenelemente über ihre Nummer ab 0):

\begin{center}
\begin{tabular}{ll}
\toprule
Zeiger & Bedeutung \\
\midrule
\code{/layers/0/thickness\_um} & Dicke der untersten Schicht \\
\code{/layers/1/breakpoints\_um/1} & zweite Lamellenkante der zweiten Schicht \\
\code{/period\_um/x} & Periode in x \\
\bottomrule
\end{tabular}
\end{center}

Bei einer Lamellenschicht ist \code{breakpoints\_um/0} die Kante bei 0 und die letzte Kante die
Periode; optimieren Sie die Kanten dazwischen. Das Plugin prüft vor dem Start, dass jeder Zeiger in
der Konfiguration steht, dort eine Zahl ist und dass der Startwert zwischen \gui{lower} und
\gui{upper} liegt; sonst meldet es den Grund (etwa \enquote{/layers/0/thickness\_um = 0.1 lies outside
its bounds [0.2, 0.5]}).

\paragraph{Weitere Einstellungen.} \gui{Iterations} (Standard 50) begrenzt die Iterationen von
L-BFGS-B. \gui{Apply the optimum to the panel} (Standard an) schreibt das Optimum nach dem Ende ins
Panel.

\paragraph{Ablauf und Ergebnisse.}
\begin{itemize}
  \item Im Fenster \gui{Plot} wächst die Kurve \enquote{objective L} über der Zahl der Auswertungen,
    live während der Rechnung. Die Statusleiste zeigt den aktuellen Wert \enquote{L = …}.
  \item \gui{Cancel} beendet die Suche nach der laufenden Auswertung; das Ergebnis ist der beste
    bis dahin gefundene Punkt.
  \item Am Ende zeigt das Fenster \gui{Table} die Tabelle \enquote{Parameters} mit den Spalten
    \code{parameter}, \code{start}, \code{optimum}, \code{lower}, \code{upper}.
  \item Das Protokoll nennt $L$ am Start und am Optimum, die Zahl der Auswertungen und die Meldung
    des Optimierers; die Statusleiste etwa \enquote{L 0.3121 $\to$ 0.6842, applied to the panel}
    (mit \enquote{(stopped)}, wenn Sie abgebrochen haben).
\end{itemize}

\begin{tipp}
Ein Gradientenverfahren findet ein \emph{lokales} Optimum in der Nähe des Starts. Probieren Sie
verschiedene Startwerte, oder verschaffen Sie sich zuerst mit einem Sweep (\gui{Scan}) einen
Überblick über die Zielfunktion. Enge Grenzen halten die Struktur herstellbar.
\end{tipp}

\begin{hinweis}
Das Optimum wird als neue Konfiguration ins Panel geschrieben, so als hätten Sie sie geladen. Das
gilt auch für einen abgebrochenen Lauf (sofern mindestens eine Auswertung stattfand). Wollen Sie den
Ausgangszustand behalten, speichern Sie die Konfiguration vorher mit \gui{Save configuration} oder
schalten \gui{Apply the optimum to the panel} aus.
\end{hinweis}

% ---------------------------------------------------------------------------------------------
\section{Nahfeld}\label{sec:fg-nahfeld}
\index{Nahfeld}

\gui{Near field} berechnet das elektromagnetische Feld im Gitter und in seiner Umgebung am ersten
Punkt der Beleuchtung und zeigt zwei Schnitte:

\begin{description}
  \item[Vertikaler Schnitt] in der x-y-Ebene, über \gui{Periods} Perioden in x und von \gui{y from}
    bis \gui{y to} in y -- als Bild im Fenster \gui{Image}. Bei einem 2D-Gitter liegt er bei
    $z$ = \gui{Vertical cut at z (2D)}.
  \item[Horizontaler Schnitt] in der Höhe \gui{Horizontal cut at y}: bei einem 1D-Gitter eine
    Kurve über x im Fenster \gui{Plot}, bei einem 2D-Gitter ein Bild in der x-z-Ebene.
\end{description}

Die Höhe $y = 0$ ist die Unterkante des Stapels (Grenze zum Substrat); negative Werte liegen im
Substrat, Werte oberhalb der Stapeldicke im Superstrat.

\subsection{Die Einstellungen}

\begin{description}
  \item[\gui{Show}] die dargestellte Größe:
    \begin{itemize}
      \item \code{|E|$^2$}, \code{|H|$^2$}: Betragsquadrat des elektrischen bzw. magnetischen
        Feldes,
      \item \code{Sy (power flow)}: der Energiefluss in y (Poynting-Vektor),
      \item \code{|Ex|} bis \code{|Hz|}: der Betrag einer Komponente,
      \item \code{Re Ex} bis \code{Re Hz}: der Realteil einer Komponente (Momentaufnahme der Welle),
      \item \code{material}: die Materialverteilung -- zur Kontrolle der Geometrie.
    \end{itemize}
    Beträge erscheinen mit der Farbskala \code{inferno}, Realteile und \code{Sy} mit
    \code{coolwarm} und einer um null symmetrischen Skala, das Material mit \code{viridis}.
  \item[\gui{y from}, \gui{y to}] Bereich des vertikalen Schnitts in y (Standard $-0{,}2$ bis
    $0{,}5$ \si{\micro\metre}).
  \item[\gui{Horizontal cut at y}] Höhe des horizontalen Schnitts (Standard 0,1 \si{\micro\metre}).
  \item[\gui{Vertical cut at z (2D)}] Lage des vertikalen Schnitts eines Kreuzgitters (Standard 0).
  \item[\gui{Points x}, \gui{Points y}, \gui{Points z (2D)}] Auflösung (Standard 400, 300, 200).
  \item[\gui{Periods}] wie viele Perioden in x (und z) gezeigt werden (Standard 2).
  \item[\gui{Also write .npz}] schreibt alle Größen beider Schnitte mit ihren Achsen in eine
    NumPy-Datei (\autoref{sec:fg-dateien}).
\end{description}

Die Normierung ist die von FourierGrating: die einfallende Welle hat $|E_\text{inc}| = 1$, das
Magnetfeld wird als $Z_0 H$ angegeben. Die Bildbeschreibung im Fenster \gui{Image} nennt sie.

\begin{hinweis}
Die Werte der Gruppe \gui{Near field} gelangen nur dann in die Konfigurationsdatei, wenn die geladene
Datei einen Block \code{nearfield} hatte oder Sie einen Wert vom Standard abweichend eingestellt
haben. Weitere Schlüssel des Blocks (\code{quantities}, \code{reconstruction}, \code{field}) stehen
unter \gui{Advanced}.
\end{hinweis}

Fehlt eine Größe im Ergebnis von FourierGrating, meldet das Plugin \enquote{the near field has no
\dots{} (add it to nearfield.quantities under Advanced)}.

\begin{tipp}
Schauen Sie sich vor einer langen Rechnung mit \gui{Show} = \code{material} den vertikalen Schnitt
an: Er zeigt sofort, ob Lamellen, Profile oder CAD-Körper so liegen, wie Sie es erwarten.
\end{tipp}

% ---------------------------------------------------------------------------------------------
\section{Schichten aus CAD-Körpern}\label{sec:fg-cad}
\index{CAD-Schicht}\index{STEP}

Komplizierte Elementarzellen -- schräge Flanken, Hinterschneidungen, mehrere Materialien --
zeichnen Sie am bequemsten im CAD. Das Plugin schreibt die Körper als STEP-Datei und legt eine
Schicht vom Typ \code{cad} an, die FourierGrating beim Rechnen in Bänder zerlegt.

\subsection{Die Körper zeichnen}

\begin{itemize}
  \item Zeichnen Sie die Körper \textbf{einer} Elementarzelle als geschlossene Körper (Schalen,
    siehe \autoref{kap:modellieren}).
  \item Die CAD-Achse $+Z$ ist die Stapelrichtung (y von FourierGrating), $+X$ die Gitterrichtung
    x. Bei einem 2D-Gitter ist die zweite Periode entsprechend die dritte Achse.
  \item Die Einheit wählen Sie mit \gui{Bodies drawn in}: \si{\micro\metre} oder \code{mm}
    (Standard Mikrometer).
  \item Jeder Körper braucht ein Material: das Attribut \code{material} des Körpers (siehe
    \autoref{kap:organisieren}) oder, wenn es fehlt, sein Name -- sofern er einem Material der
    Materialliste entspricht. Ein Körper ohne passendes Material wird mit einer Meldung abgewiesen,
    die die verfügbaren Materialien nennt.
\end{itemize}

\subsection{Die Schicht anlegen}

\begin{schritte}
  \item Wählen Sie die Körper im Ansichtsfenster aus und drücken Sie in der Gruppe \gui{CAD layer}
    bei \gui{Bodies} auf \gui{Use selection}. Das Feld zeigt die Namen der Körper.
  \item Wählen Sie unter \gui{STEP file} die Datei, in die die Körper geschrieben werden. Die Schicht
    verweist mit ihrem absoluten Pfad auf diese Datei.
  \item Wählen Sie \gui{Around the bodies}: das Material zwischen den Körpern, \code{superstrate}
    (Standard) oder \code{substrate}.
  \item Stellen Sie bei \gui{RCWA slices per band} ein, in wie viele Scheiben jedes Band der
    RCWA zerlegt wird (Standard 30).
  \item Wählen Sie \menu{Plugins,FourierGrating,Layer from CAD bodies}.
\end{schritte}

Das Plugin schreibt die Körper (in der gewählten Reihenfolge) als STEP, hängt eine Schicht vom Typ
\code{json} mit dem Eintrag \code{\textquotedbl{}type\textquotedbl{}: \textquotedbl{}cad\textquotedbl{}} \textbf{oben} an den Stapel an und prüft sie
sofort (wie \gui{Check CAD layers}). Die Statusleiste meldet etwa \enquote{3 bodies in
zelle.step, cad layer 2 added (top); 4 bands, 0 errors, 0 warnings}. Mit $\uparrow$/$\downarrow$
schieben Sie die Schicht an ihren Platz im Stapel.

Hat das installierte \code{fouriergrating} keinen CAD-Slicer, wird die Schicht trotzdem angelegt;
das Protokoll meldet \enquote{not checked: this fouriergrating has no CAD slicer …}.

\subsection{Prüfen: Check CAD layers}

\gui{Check CAD layers} führt die Geometrieprüfung von FourierGrating für jede \code{cad}-Schicht des
Stapels aus und zeigt im Fenster \gui{Table}

\begin{description}
  \item[\enquote{Bands}] die Zerlegung in Bänder: Schicht, \code{y from / }\si{\micro\metre}, \code{y to / }\si{\micro\metre},
    Art des Bandes (homogen, Extrusion, geschnitten) und die Flächenanteile je Material;
  \item[\enquote{Findings}] (nur wenn es welche gibt) die Befunde: Schicht, Schwere (\code{error}
    oder \code{warning}), Code, Meldung und Ort in \si{\micro\metre}.
\end{description}

Die Statusleiste fasst zusammen: \enquote{N bands, E errors, W warnings}, mit dem Zusatz
\enquote{– not usable}, wenn die Schicht so nicht gerechnet werden kann.

Steht unter \gui{Preview SVG} eine Datei, schreibt die Prüfung dorthin zusätzlich den Schnitt der
ersten \code{cad}-Schicht in der Höhe \gui{Preview at height} -- ein Bild, das Sie in jedem Browser
öffnen können.

% ---------------------------------------------------------------------------------------------
\section{Konfigurationsdateien}\label{sec:fg-dateien}
\index{Konfigurationsdatei (FourierGrating)}

\subsection{Laden und Speichern}

FourierGrating speichert ein Gitter als JSON-Datei. Das Plugin liest und schreibt dieses Format so,
dass Laden und erneutes Speichern dieselbe Datei ergibt: auch die Reihenfolge der Schlüssel bleibt
erhalten.

\begin{description}
  \item[\gui{Load configuration}] liest die Datei aus \gui{Configuration} (Gruppe \gui{File}) ins
    Panel. Struktur, Beleuchtung, Materialien, Schichten und Nahfeldeinstellungen werden ersetzt;
    die Einstellungen von Scan, Optimierung und CAD-Schicht bleiben. Die Statusleiste nennt die
    Zahl der Materialien und Schichten (\enquote{(some as JSON)}, wenn Schichten als JSON
    erscheinen).
  \item[\gui{Save configuration}] schreibt die Konfiguration des Panels in die Datei aus
    \gui{Save as} -- oder, wenn dieses Feld leer ist, in die Datei aus \gui{Configuration}.
\end{description}

\begin{warnung}
Ist \gui{Save as} leer, überschreibt \gui{Save configuration} die geladene Datei ohne Rückfrage.
Wählen Sie unter \gui{Save as} eine neue Datei, wenn Sie das Original behalten wollen.
\end{warnung}

\subsection{Weitere Einstellungen (Advanced)}\label{sec:fg-advanced}

Das Panel zeigt nur die Teile einer Konfiguration als Felder, die man am häufigsten bearbeitet.
Alles andere -- Solver-Optionen (\code{rcwa}, \code{rk}, \code{backend}, \code{precision},
\code{threads}, \dots), ein \code{sweep}, die Ausgabe (\code{output}), weitere Schlüssel von
\code{incidence}, \code{truncation} und \code{nearfield} -- steht unverändert als JSON-Objekt in
\gui{Further settings (JSON)} der Gruppe \gui{Advanced}. Beim Speichern wird es wieder in die
Konfiguration eingefügt. Die Bedeutung dieser Schlüssel beschreibt die Formatdokumentation von
FourierGrating (\datei{docs/config\_format.md} dort).

Ein Beispiel mit mehr RCWA-Scheiben und höchstens vier parallelen Scanpunkten:

\begin{lstlisting}
{ "rcwa": { "slices": 60 }, "threads": { "scan": 4 } }
\end{lstlisting}

\begin{hinweis}
Schlüssel wie \code{dimension}, \code{solver}, \code{truncation.Nx} oder \code{period\_um.z}, die die
geladene Datei nicht hatte, schreibt das Plugin nur, wenn Sie sie vom Standardwert abweichend
eingestellt haben. So bleibt eine Datei, die Sie nur laden und wieder speichern, unverändert.
\end{hinweis}

\subsection{Exportierte Dateien}

\begin{referenz}{>{\raggedright\arraybackslash}p{3.6cm}>{\raggedright\arraybackslash}p{3.2cm}>{\raggedright\arraybackslash}X}
\toprule
Feld & Aktion & Inhalt \\
\midrule
\endhead
\gui{Save as} & \gui{Save configuration} & Konfiguration von FourierGrating (JSON, eingerückt) \\
\gui{Solve: also write CSV} & \gui{Solve} & Ordnungstabelle: \code{side, n, m, efficiency, eff\_s, eff\_p} \\
\gui{Also write CSV} (\gui{Scan}) & \gui{Scan} & jede Ordnung an jedem Punkt wie \datei{efficiencies.csv} \\
\gui{Also write .npz} & \gui{Near field} & alle Größen beider Schnitte mit Achsen und Metadaten \\
\gui{STEP file} & \gui{Layer from CAD bodies} & die gewählten Körper als STEP \\
\gui{Preview SVG} & \gui{Check CAD layers} & Schnitt der ersten \code{cad}-Schicht \\
\bottomrule
\end{referenz}

Die genauen Spalten und Feldnamen stehen im Anhang (\autoref{kap:plugin-fouriergrating}).

% ---------------------------------------------------------------------------------------------
\section{Das Gitter als Block im Workflow}\label{sec:fg-block}
\index{Gitterblock}\index{Jones-Tabelle}

Im Node-Workflow (\autoref{kap:workflow}) beschreiben Blöcke einen optischen Aufbau -- Quellen,
Linsen, Spiegel, Kameras. Das FourierGrating-Plugin liefert zwei Blocktypen, mit denen ein
rigoros gerechnetes Gitter oder ein Schichtstapel Teil eines solchen Aufbaus wird:

\begin{description}
  \item[\block{fg.grating}] eine periodische Probe, die jede gewünschte Beugungsordnung als eigenen
    Kanal weitergibt -- ein Strahlteiler, dessen Ausgänge die Ordnungen sind. Er hat einen Eingang
    \code{in} und einen Ausgang \code{out} (\enquote{orders}) und nimmt den Platz eines Spiegels
    oder einer Platte ein. In älteren Dokumenten hieß der Typ \block{fusion.grating}; solche Dokumente
    werden weiterhin geladen.
  \item[\block{fg.stack}] ein lateral unendlich ausgedehnter Stapel homogener Schichten zwischen zwei
    Medien -- Beschichtung, Filter, Spiegelschichten --, exakt für jede Strahlrichtung gerechnet
    (FourierGrating mit $N = 0$), mit Licht von beiden Seiten (Ein- und Ausgänge \code{front\_in},
    \code{front\_out}, \code{back\_in}, \code{back\_out}).
\end{description}

Die Workflow-Engine schickt die Auswertung dieser Blöcke an den Prozess des Plugins
FourierGrating.
Ports, Parameter und ihre Bedeutung stehen im Anhang (\autoref{sec:fg-ref-bloecke}) und in
\autoref{kap:workflow-bloecke}.

\subsection{Jones-Tabelle oder Konfiguration}

Ein Gitterblock kennt zwei Arten, die Wirkung des Gitters zu beschreiben:

\begin{itemize}
  \item eine \textbf{Jones-Tabelle}: für jede Ordnung die leistungsnormierten Jones-Matrizen über
    Wellenlänge und Einfallsrichtung, vorab berechnet (Parameter \param{table}); der Block braucht
    dann keinen FourierGrating-Rechenkern;
  \item eine \textbf{FourierGrating-Konfiguration}: der Block baut die Tabelle selbst, passend zu
    den einfallenden Strahlen -- um deren Wellenlänge und den Bereich ihrer Richtungskosinus
    $\sigma$, mit einem Rand und Stützstellen im Abstand \gui{$\sigma$ step of the table} (Gruppe
    \gui{Ray path}). Für ein Feld (Darstellung H) rechnet er jede ebene Welle des Spektrums.
\end{itemize}

Hält ein Block eine Jones-Tabelle, wird diese verwendet und nicht die Konfiguration.

\subsection{Die Konfiguration liegt im Dokument}

Seit RTT-Fusion W75 hält ein Gitter- oder Stapelblock seine Konfiguration selbst
(Parameter \param{config}). Ein \datei{.rttf}-Dokument rechnet damit auf jedem Rechner, ohne weitere
Dateien. Ältere Blöcke, die noch auf eine Datei verweisen (Parameter \param{fouriergrating}),
bekommen deren Inhalt, sobald \gui{Check workflow} oder \gui{Run workflow} die Datei lesen kann; der
Pfad wird dann geleert.

Zwischen Panel und Block wechseln Sie mit zwei Aktionen im \textbf{Kontextmenü des Blocks} im
Node-Editor (\maus{Rechtsklick} auf den Block):

\begin{description}
  \item[\gui{Store panel configuration in block}] Die Konfiguration des Panels wird die des Blocks.
    Der Block verweist danach auf keine Datei mehr; eine Jones-Tabelle eines Gitterblocks wird
    verworfen (sie würde sonst statt der Konfiguration verwendet). Die Statusleiste meldet die Zahl
    der Materialien und Schichten und, falls zutreffend, \enquote{its Jones table is dropped}.
  \item[\gui{Load block configuration into panel}] zeigt die Konfiguration des Blocks im Panel. Dort
    können Sie sie bearbeiten, mit \gui{Solve}, \gui{Scan} und \gui{Optimise} untersuchen und mit
    \gui{Store panel configuration in block} zurückschreiben. Das Feld \gui{Configuration} wird
    dabei geleert, damit \gui{Save configuration} nicht versehentlich eine fremde Datei überschreibt.
\end{description}

Ein Block, der nur eine Jones-Tabelle hat, besitzt keine Konfiguration; \gui{Load block
configuration into panel} meldet dann \enquote{… holds no FourierGrating configuration}. Beide
Aktionen funktionieren auch für Blöcke in Gruppen (in der aktiven Variante).

\subsection{Der Dialog des Blocks}

Ein \maus{Doppelklick} auf den Block öffnet seinen Dialog. Hat der Block eine Konfiguration, zeigt
der Dialog deren wichtigste Zahlen als Felder, die Sie direkt ändern können -- gleich, ob die
Konfiguration aus dem Panel, aus einer eingebetteten Datei oder mit dem Dokument kam (das Plugin
beschreibt seine Blöcke nach jeder Änderung des Workflows, beim Öffnen eines Dokuments und wenn es
startet; es muss dazu laufen):

\begin{description}
  \item[\gui{Grating}] \gui{Period x} (in 2D auch \gui{Period z}) und
    \gui{Fourier orders Nx ($-$N … N)} (in 2D auch \gui{Nz});
  \item[\gui{Layers}] je Schicht \gui{Layer $i$ (Typ): thickness} und bei Lamellen
    \gui{Layer $i$ (Typ): lamella edges} -- die Schichten von 1 an gezählt;
  \item[\gui{Materials}] je Material \gui{Name: n} oder \gui{Name: n, k}.
\end{description}

Daneben stehen die Einstellungen des Blocks selbst: Lage (\param{point}, \param{normal}, \param{grating\_\allowbreak vector}), die weitergegebene Seite (\param{side}: \code{reflected} oder
\code{transmitted}), die Auswahl der Ordnungen (\param{orders}, leer: alle), die Tabellenauflösung
(\gui{Ray path}) und die Einstellungen des Feldpfads (\gui{Field path}).

\begin{tipp}
Für einen Stapelblock speichern Sie im Panel eine Konfiguration mit ausschließlich homogenen
Schichten, von unten (Substrat, Rückseite des Blocks) nach oben (Superstrat, Vorderseite, in
Richtung \param{normal}). Eine Lamellen- oder Profilschicht weist der Stapelblock ab: sie gehört in
einen \block{fg.grating}.
\end{tipp}

% ---------------------------------------------------------------------------------------------
\section{Arbeitsabläufe Schritt für Schritt}\label{sec:fg-ablaeufe}

\subsection{Die erste Rechnung: eine Glasoberfläche}\label{sec:fg-ablauf-erste}

Diese Rechnung prüft, ob das Plugin und seine Umgebung funktionieren, und hat ein bekanntes
Ergebnis: die Fresnel-Reflexion von Glas.

\begin{schritte}
  \item Öffnen Sie das Dock \gui{Plugin} und wählen Sie oben \gui{FourierGrating}.
  \item Drücken Sie \gui{Reset parameters}. Das Panel enthält jetzt die Standardwerte: Luft über
    Glas ($n = 1{,}5$), keine Schichten, $\lambda = 0{,}6328$ \si{\micro\metre}, senkrechter Einfall.
  \item Drücken Sie \gui{Solve}.
  \item Im Fenster \gui{Table} steht die Ordnungstabelle. Die Ordnung \code{R 0} hat die Effizienz
    $((1{,}5-1)/(1{,}5+1))^2 = 0{,}04$, \code{T 0} den Rest $0{,}96$; die übrigen Ordnungen sind
    ohne Gitter null.
  \item Im \gui{Plugin log} stehen $\Sigma R = 0{,}04$, $\Sigma T = 0{,}96$ und eine Absorption $A$ von (nahezu) null.
\end{schritte}

Erscheint stattdessen die Meldung über das fehlende Paket \code{fouriergrating}, richten Sie die
Umgebung ein (\autoref{sec:fg-voraussetzungen}).

\subsection{Ein binäres Reflexionsgitter}\label{sec:fg-ablauf-lamellen}

Ein Lamellengitter aus Glas auf Glas, \qty{0.3}{\micro\metre} tief, Periode
\qty{1}{\micro\metre}, Tastverhältnis 0,5.

\begin{schritte}
  \item Setzen Sie in \gui{Structure} \gui{Period x} = 1 und lassen Sie \gui{Orders Nx} = 10.
  \item Fügen Sie unter \gui{Layers (bottom to top)} mit \gui{Add} eine Schicht hinzu und öffnen Sie
    sie mit \gui{Edit…}.
  \item Wählen Sie \gui{Type} = \code{lamellar1d}, \gui{Thickness} = 0.3.
  \item Tragen Sie in \gui{Lamellae (lamellar1d)} zwei Zeilen ein: \code{0.0 glass} und
    \code{0.5 air}. Setzen Sie \gui{End (lamellar1d)} = 1.0. Schließen Sie den Dialog.
  \item Drücken Sie \gui{Solve}. Die Tabelle zeigt jetzt Effizienzen in den Ordnungen $\pm 1$ und
    höher; $\Sigma R + \Sigma T = 1$ (verlustfrei).
  \item Prüfen Sie die Konvergenz: setzen Sie \gui{Orders Nx} = 20 und rechnen Sie erneut. Die
    Effizienzen sollten sich nur noch in den hinteren Stellen ändern.
\end{schritte}

\subsection{Spektrum einer Ordnung in beiden Polarisationen}

\begin{schritte}
  \item Ausgehend vom Gitter des vorigen Abschnitts: tragen Sie bei \gui{$\lambda$ /
    \si{\micro\metre}} den Bereich \code{0.4:0.8:41} ein und bei \gui{$\delta$ / deg} die Liste
    \code{0, 90}.
  \item Drücken Sie \gui{Estimate}: 82 Rechnungen, Matrixgröße 84.
  \item Lassen Sie \gui{x axis} auf \code{auto} (die Wellenlänge hat die meisten Werte) und drücken
    Sie \gui{Scan}.
  \item Das Fenster \gui{Plot} zeigt $\Sigma R$, $\Sigma T$ und alle Ordnungen über 1\,\% -- jede
    zweimal, mit \code{($\delta$ = 0)} und \code{($\delta$ = 90)}.
  \item Wollen Sie die Zahlen weiterverarbeiten, wählen Sie unter \gui{Also write CSV} eine Datei
    und scannen erneut.
\end{schritte}

Die Ordnungen $\pm 1$ verschwinden im Superstrat, sobald $\lambda > \Lambda_x$ (bei senkrechtem
Einfall in Luft): ihre Kurven enden dort.

\subsection{Ein Littrow-Gitter optimieren}\label{sec:fg-ablauf-optimieren}

Ziel: möglichst viel Licht in die $-1$. reflektierte Ordnung unter Littrow-Einfall, durch Wahl von
Tiefe und Lamellenbreite.

\begin{schritte}
  \item Bauen Sie das Lamellengitter aus \autoref{sec:fg-ablauf-lamellen} auf, eine Wellenlänge
    (\code{0.6328}), \gui{$\delta$} = \code{90}.
  \item Setzen Sie \gui{Littrow order} = $-1$ und wählen Sie \gui{Set Littrow angle}. $\varphi$
    wird \ang{18.4454}.
  \item Rechnen Sie einmal \gui{Solve} und notieren Sie die Effizienz von \code{R -1}.
  \item Tragen Sie in \gui{Objective L = $\Sigma$ weight $\cdot$ efficiency} eine Zeile ein:
    \gui{side} \code{R}, \gui{n} \code{-1}, \gui{m} \code{0}, \gui{weight} \code{1}. \gui{Goal}
    bleibt \code{maximise}.
  \item Tragen Sie in \gui{Parameters} zwei Zeilen ein:
    \code{/layers/0/thickness\_um} mit den Grenzen 0.05 und 0.8, und
    \code{/layers/0/breakpoints\_um/1} mit den Grenzen 0.2 und 0.8.
  \item Drücken Sie \gui{Optimise}. Verfolgen Sie im \gui{Plot} die Kurve \enquote{objective L}.
  \item Nach dem Ende stehen die optimierten Werte in der Tabelle \enquote{Parameters} und -- mit
    \gui{Apply the optimum to the panel} -- im Panel. \gui{Solve} bestätigt das Ergebnis.
\end{schritte}

\begin{hinweis}
Ändern Sie mit der Optimierung eine Lamellenkante, bleibt das Ende (\gui{End}) bei der Periode.
Optimieren Sie die Periode selbst (\code{/period\_um/x}), passt das Plugin die Lamellenkanten
\emph{nicht} an -- die letzte Kante wäre danach falsch (\autoref{sec:fg-lamellen}). Optimieren Sie
die Periode nur bei Schichten ohne Lamellen oder nehmen Sie die letzte Kante als weiteren Parameter
mit gleichen Grenzen auf.
\end{hinweis}

\subsection{Eine Schichtdicke durchfahren (Sweep)}

\begin{schritte}
  \item Setzen Sie alle Beleuchtungsfelder auf einzelne Zahlen.
  \item Tragen Sie unter \gui{Advanced} ein:

\begin{lstlisting}
{ "sweep": { "parameter": "/layers/0/thickness_um",
             "values": { "from": 0.05, "to": 0.5, "steps": 46 } } }
\end{lstlisting}
  \item Wählen Sie \gui{x axis} = \code{sweep\_value} (oder \code{auto}) und drücken Sie \gui{Scan}.
  \item Die Kurven zeigen die Effizienzen über der Schichtdicke; die x-Achse trägt den Namen des
    Parameters.
\end{schritte}

Entfernen Sie den Sweep danach wieder aus \gui{Advanced}, sonst rechnet jeder weitere Scan ihn mit.

\subsection{Das Nahfeld eines Gitters ansehen}

\begin{schritte}
  \item Stellen Sie das Gitter ein und prüfen Sie die Effizienzen mit \gui{Solve}.
  \item Wählen Sie in der Gruppe \gui{Near field} \gui{Show} = \code{material} und drücken Sie
    \gui{Near field}. Das Bild zeigt die Geometrie über zwei Perioden.
  \item Passen Sie \gui{y from}/\gui{y to} so an, dass der Stapel und etwas Superstrat und Substrat
    sichtbar sind -- bei \qty{0.3}{\micro\metre} Tiefe etwa $-0{,}2$ bis $0{,}6$.
  \item Wählen Sie \gui{Show} = \code{|E|$^2$} und rechnen Sie erneut. Für eine Momentaufnahme der
    Welle nehmen Sie \code{Re Ez} (s-Polarisation in 1D) oder \code{Re Hz} (p-Polarisation).
  \item Für die Auswertung in Python wählen Sie unter \gui{Also write .npz} eine Datei.
\end{schritte}

\subsection{Eine Elementarzelle aus dem CAD rechnen}

\begin{schritte}
  \item Legen Sie im Panel die Materialien an, z.\,B. \code{si} mit $n = 3{,}88$, $k = 0{,}02$.
  \item Zeichnen Sie die Körper einer Elementarzelle (Einheit Mikrometer, Stapelrichtung $+Z$,
    Gitterrichtung $+X$), z.\,B. einen Quader von $x = 0$ bis $0{,}4$ und $z = 0$ bis $0{,}3$.
  \item Geben Sie jedem Körper das Attribut \code{material} mit dem Materialnamen -- oder benennen
    Sie den Körper wie das Material.
  \item Wählen Sie die Körper aus, drücken Sie bei \gui{Bodies} auf \gui{Use selection} und wählen
    Sie eine \gui{STEP file}.
  \item Drücken Sie \gui{Layer from CAD bodies}. Die Tabelle \enquote{Bands} zeigt die Zerlegung;
    eine Tabelle \enquote{Findings} erscheint nur bei Problemen.
  \item Setzen Sie die Periode passend zur gezeichneten Zelle und rechnen Sie mit \gui{Solve}.
\end{schritte}

\subsection{Eine Konfiguration in einen Gitterblock übernehmen}

So bringen Sie ein Gitter, das Sie in GUI-FourierGrating oder im Panel aufgebaut haben, in den
Node-Workflow:

\begin{schritte}
  \item Laden Sie die Konfiguration: Datei unter \gui{Configuration} wählen, \gui{Load configuration}
    -- oder bauen Sie sie im Panel auf.
  \item Legen Sie im Node-Editor einen Block \block{fg.grating} an (oder ersetzen Sie einen Spiegel
    mit \gui{Replace with}), siehe \autoref{kap:workflow}.
  \item \maus{Rechtsklick} auf den Block, \gui{Store panel configuration in block}.
  \item Stellen Sie im Dialog des Blocks (\maus{Doppelklick}) Lage, Gittervektor, Seite und Ordnungen
    ein.
  \item Prüfen und rechnen Sie den Workflow mit \gui{Check workflow} und \gui{Run workflow}.
\end{schritte}

Später ändern: \maus{Rechtsklick}, \gui{Load block configuration into panel}, bearbeiten, wieder
\gui{Store panel configuration in block}.

% ---------------------------------------------------------------------------------------------
\section{Typische Fehler und Meldungen}\label{sec:fg-fehler}

\begin{referenz}{>{\raggedright\arraybackslash}p{6.2cm}>{\raggedright\arraybackslash}X}
\toprule
Meldung & Ursache und Abhilfe \\
\midrule
\endhead
\code{The FourierGrating plugin needs the Python package fouriergrating …} & Paket fehlt in der
  Plugin-Umgebung: \gui{Create environment…} oder \code{pip install fouriergrating}. \\
\code{Choose the configuration first: File $\to$ Configuration.} & \gui{Load configuration} ohne Datei:
  unter \gui{Configuration} eine Datei wählen. \\
\code{Choose where to save: File $\to$ Save as.} & \gui{Save configuration} ohne Ziel. \\
\code{The panel does not give a valid configuration: …} & Ein Feld der Beleuchtung lässt sich nicht
  lesen (z.\,B. Komma statt Punkt, \code{0,5:0,7:21}), oder eine Tabelle ist unvollständig. \\
\code{FourierGrating: …} & Meldung des Rechenkerns, etwa ein unbekannter Materialname, ein
  verlustbehaftetes Superstrat oder eine fehlerhafte JSON-Schicht. Der Text nennt meist den
  JSON-Zeiger. \\
\code{the scan has a single point …} & \gui{Scan} ohne Liste, Bereich oder Sweep. \\
\code{… does not vary in this scan (varying: …)} & Die gewählte \gui{x axis} variiert nicht. \\
\code{no objective terms – add a row (side, n, m, weight)} & Tabelle der Zielfunktion leer. \\
\code{no parameters – add a row (JSON pointer, lower, upper)} & Parametertabelle leer. \\
\code{… is not in the configuration} & JSON-Zeiger falsch, z.\,B. Schicht nicht vorhanden oder
  Tippfehler im Schlüssel. \\
\code{… lies outside its bounds […]} & Startwert außerhalb der Grenzen: Grenzen erweitern oder Wert
  im Panel ändern. \\
\code{Littrow needs a period x > 0} & \gui{Period x} ist 0. \\
\code{order … has no Littrow angle here …} & Ordnung 0 oder $|\sin\varphi| > 1$: andere Ordnung oder
  kürzere Wellenlänge. \\
\code{Choose the bodies first: …} & \gui{Layer from CAD bodies} ohne Körper. \\
\code{body … has no material …} & Körper ohne Attribut \code{material} und ohne passenden Namen. \\
\code{No cad layer – add one with Layer from CAD bodies (or as JSON).} & \gui{Check CAD layers}
  ohne \code{cad}-Schicht. \\
\code{this fouriergrating has no CAD slicer …} & \code{fouriergrating} ohne
  \code{FG\_ENABLE\_NURBFORGE}. \\
\code{The document has no node workflow.} / \code{The block … is gone …} & Blockaktion, während der
  Workflow fehlt oder der Block gelöscht wurde. \\
\bottomrule
\end{referenz}

\begin{warnung}
Die häufigste stille Fehlerquelle sind Lamellenkanten, die nicht bei der Periode enden
(\autoref{sec:fg-lamellen}). Prüfen Sie nach jeder Änderung der Periode das Feld \gui{End} -- oder
sehen Sie sich mit \gui{Near field} und \gui{Show} = \code{material} die Geometrie an.
\end{warnung}

% ---------------------------------------------------------------------------------------------
\section{Grenzen}\label{sec:fg-grenzen}

\begin{itemize}
  \item \gui{Solve}, \gui{Near field} und \gui{Optimise} rechnen nur den ersten Punkt einer
    Beleuchtungsliste; das Nahfeld eines Sweeps gilt für dessen ersten Wert.
  \item Eine laufende Rechnung von FourierGrating lässt sich nicht unterbrechen
    (\autoref{sec:fg-ueberblick}); große Rechnungen vorher abschätzen.
  \item Die Speicherschätzung ist eine Faustregel und kann deutlich danebenliegen.
  \item Das Panel zeigt nur homogene und Lamellenschichten sowie Materialien mit Brechzahl als
    Felder; alles andere bearbeiten Sie als JSON. Das Plugin prüft JSON-Einträge nicht selbst.
  \item Lamellenkanten werden nicht gegen die Periode geprüft.
  \item Die Optimierung findet lokale Optima; die Parameter müssen Zahlen der Konfiguration sein.
  \item CAD-Schichten brauchen ein \code{fouriergrating} mit CAD-Slicer. Die Schicht verweist auf die
    STEP-Datei über ihren absoluten Pfad.
  \item Die Workflow-Blöcke gibt es nur mit einem installierten \code{rtt-fusion}, das Fernblöcke
    unterstützt; ohne das Plugin rechnet die Workflow-Engine einen Gitterblock selbst, mit demselben
    Ergebnis.
\end{itemize}
